% bismi ALLAH ARRAHMAN ARRAHIM

\documentclass[11pt,a4paper]{article}

\usepackage[utf8]{inputenc}
\usepackage[T1]{fontenc}
\usepackage[french]{babel}
\usepackage{amsmath,amssymb,amsfonts,amsthm,mathtools}
\usepackage{booktabs}
\usepackage{geometry}
\usepackage{enumitem}
\usepackage{microtype}
\usepackage{hyperref}
\hypersetup{hidelinks}
\usepackage{tikz}
\usetikzlibrary{arrows.meta,positioning,calc}

\geometry{margin=2.5cm}

% -------------------------------------------------
% Environnements
% -------------------------------------------------
\newtheorem{definition}{Définition}[section]
\newtheorem{assumption}{Hypothèse}[section]
\newtheorem{proposition}{Proposition}[section]
\newtheorem{theorem}{Théorème}[section]
\newtheorem{corollary}{Corollaire}[section]
\newtheorem{remark}{Remarque}[section]
\newtheorem{lemma}{Lemme}[section]
\newtheorem{example}{Exemple}[section]

% -------------------------------------------------
% Opérateurs
% -------------------------------------------------
\DeclareMathOperator{\prox}{prox}
\DeclareMathOperator{\argmin}{argmin}
\DeclareMathOperator{\argmax}{argmax}
\DeclareMathOperator{\dist}{dist}
\DeclareMathOperator{\conv}{conv}

% -------------------------------------------------
% Raccourcis
% -------------------------------------------------
\newcommand{\R}{\mathbb R}
\newcommand{\Hh}{\mathcal H}
\newcommand{\Mcal}{\mathcal M}
\newcommand{\Ccal}{\mathcal C}
\newcommand{\Pcal}{\mathcal P}
\newcommand{\Tcal}{\mathcal T}
\newcommand{\Rcal}{\mathcal R}
\newcommand{\Acal}{\mathcal A}
\newcommand{\Bcal}{\mathcal B}

\title{Routage variationnel PCH-ESM\\
\large Sélection de contexte par fonctions valeur}
\author{}
\date{}

\begin{document}

\maketitle

\begin{abstract}
Ce manuscrit formalise le routage dans le cadre PCH comme une conséquence directe des valeurs locales déjà présentes dans le noyau de purification. Le routage n'est pas une couche ajoutée : il compare des valeurs réduites issues de problèmes locaux, puis sélectionne le contexte minimisant l'enveloppe inférieure correspondante. On obtient ainsi un problème d'optimisation sur \(\R^d\times\Ccal\), une relaxation entropique sur le simplexe des contextes, et une géométrie de décision composée de cellules, de frontières et de régularité au sens de Clarke. Les outils de Moreau, Berge, Clarke et Rockafellar--Wets sont employés uniquement là où ils sont nécessaires. Une compatibilité explicite avec la décomposition spectrale de PCH-ESM est également donnée.
\end{abstract}

\tableofcontents

% =================================================
\section{Introduction et motivation}
% =================================================

On considère un noyau PCH déjà défini et validé. Pour chaque entrée latente \(w\in\R^d\), plusieurs contextes admissibles \(c\in\Ccal\) sont possibles, chacun associé à un ensemble convexe non vide, fermé, \(\Mcal_c\subset\R^d\), et à une fonctionnelle locale propre.

Le point essentiel est le suivant : la question \emph{quel contexte purifier ?} n'est pas extérieure à PCH. Elle est déjà contenue dans la famille des valeurs locales. Le routage n'ajoute donc pas une couche artificielle ; il explicite la comparaison des branches d'énergie que PCH fournit déjà.

L'originalité du manuscrit est précisément là. Dans une mixture of experts classique, un gate appris sélectionne un expert. Ici, le contexte est sélectionné comme minimiseur d'une famille de fonctions valeur construites à partir de problèmes variationnels locaux. Le routage devient alors un objet d'analyse à part entière.

Cette lecture est aussi compatible avec l'architecture PCH-ESM : une décomposition spectrale fournit une famille finie de blocs spécialisés, et le routage peut alors choisir le bloc dont la valeur réduite est la plus basse. PCH n'est donc pas remplacé ; il est prolongé.

\begin{figure}[ht]
\centering
\begin{tikzpicture}[node distance=1.65cm, >=Latex, every node/.style={align=center}]
\node (w) [draw, rounded corners, minimum width=2.8cm, minimum height=0.9cm] {$w\in\R^d$};
\node (locals) [draw, rounded corners, right=2.7cm of w, minimum width=4.1cm, minimum height=1.0cm] {Problèmes locaux\\$\widetilde J_c(\cdot;w)$};
\node (values) [draw, rounded corners, right=2.7cm of locals, minimum width=3.8cm, minimum height=1.0cm] {Valeurs réduites\\$\widetilde V_c(w)$};
\node (argmin) [draw, rounded corners, below=1.25cm of values, minimum width=3.8cm, minimum height=1.0cm] {Sélection\\$\argmin_{c\in\Ccal}\widetilde V_c(w)$};
\node (chosen) [draw, rounded corners, left=2.7cm of argmin, minimum width=3.8cm, minimum height=1.0cm] {Contexte choisi\\$c^\star(w)$};

\draw[->] (w) -- (locals);
\draw[->] (locals) -- (values);
\draw[->] (values) -- (argmin);
\draw[->] (argmin) -- (chosen);
\end{tikzpicture}
\caption{Chaîne conceptuelle du routage variationnel.}
\end{figure}

\section*{Table des notations principales}
\begin{center}
\begin{tabular}{@{}ll@{}}
\toprule
Notation & Signification \\ \midrule
\(\widetilde J_c(u;w)\) & énergie locale normalisée du contexte \(c\) \\
\(\widetilde V_c(w)\) & valeur réduite associée au contexte \(c\) \\
\(u_c^\star(w)\) & minimiseur local unique du problème normalisé \\
\(F(w)\) & valeur globale de routage, \(F(w)=\min_c \widetilde V_c(w)\) \\
\(\Acal(w)\) & ensemble des contextes actifs \\
\(\Omega_c^{\mathrm{str}}\) & cellule stricte du contexte \(c\) \\
\(\Bcal\) & ensemble des points de bord \\
\(\Rcal_\beta\) & fonctionnelle de routage souple \\
\(\Delta_K\) & simplexe des contextes \\ \bottomrule
\end{tabular}
\end{center}

% =================================================
\section{Problème local PCH et normalisation canonique}
% =================================================

\begin{assumption}[Espace latent]
On fixe un espace euclidien réel \(\R^d\). Les états latents purifiés vivent dans \(\R^d\).
\end{assumption}

\begin{assumption}[Famille de contextes]
On considère un ensemble fini de contextes admissibles
\[
\Ccal=\{1,\dots,K\}.
\]
À chaque contexte \(c\in\Ccal\) est associé un ensemble convexe non vide, fermé, \(\Mcal_c\subset\R^d\).
\end{assumption}

\begin{assumption}[Paramètres PCH]
Pour chaque contexte \(c\), on fixe des paramètres \(\lambda_{\mathrm{ctx},c}>0\) et \(\lambda_{\mathrm{ridge},c}\ge 0\).
\end{assumption}

\begin{definition}[Énergie contextuelle]
Pour un contexte \(c\in\Ccal\), on définit l'énergie
\[
J_c(u;w)=\frac12\|u-w\|^2+\lambda_{\mathrm{ctx},c}\,\dist(u,\Mcal_c)^2+\lambda_{\mathrm{ridge},c}\,\|u\|^2.
\]
\end{definition}

\begin{definition}[Opérateur PCH associé]
Pour un contexte \(c\in\Ccal\), on définit
\[
\Pcal_c(w):=\argmin_{u\in\R^d}J_c(u;w).
\]
\end{definition}

\begin{proposition}[Convexité forte et unicité]
Pour tout contexte \(c\) et toute entrée \(w\), la fonctionnelle \(u\mapsto J_c(u;w)\) est \(\mu_c\)-fortement convexe, où
\[
\mu_c:=1+2\lambda_{\mathrm{ridge},c}.
\]
En particulier, \(\Pcal_c(w)\) est un singleton.
\end{proposition}

\begin{proof}
Le terme \(u\mapsto \frac12\|u-w\|^2+\lambda_{\mathrm{ridge},c}\|u\|^2\) est \((1+2\lambda_{\mathrm{ridge},c})\)-fortement convexe. La fonction \(u\mapsto \dist(u,\Mcal_c)^2\) est convexe parce que \(\Mcal_c\) est non vide, fermé et convexe. La somme est donc fortement convexe, d'où l'existence et l'unicité du minimiseur.
\end{proof}

\begin{remark}
Le noyau PCH n'est pas modifié ici : on ne change ni la forme de \(J_c\), ni la définition des ensembles \(\Mcal_c\), ni la logique de purification. Le routage est ajouté comme problème de sélection séparé.
\end{remark}

\begin{definition}[Fonction de base normalisée]
On pose
\[
f_c(u):=\lambda_{\mathrm{ctx},c}\,\dist(u,\Mcal_c)^2+\lambda_{\mathrm{ridge},c}\,\|u\|^2,
\qquad
g_c(u):=\frac{1}{\mu_c}f_c(u).
\]
Alors
\[
\widetilde J_c(u;w)=\frac{1}{\mu_c}J_c(u;w)=g_c(u)+\frac{1}{2\mu_c}\|u-w\|^2.
\]
\end{definition}

\begin{proposition}[Pourquoi normaliser]
La normalisation par \(\mu_c\) compare les branches sur une échelle commune et préserve le minimiseur de chaque problème local. Elle isole aussi la part quadratique du problème, ce qui permet de lire \(\widetilde J_c\) comme une enveloppe de Moreau de paramètre \(\mu_c\). En particulier, la comparaison des valeurs réduites n'est pas biaisée par une rigidité locale plus forte.
\end{proposition}

\begin{proof}
Le point de préservation du minimiseur est immédiat : multiplier par un scalaire strictement positif ne change pas l'ensemble des minimiseurs. L'identité
\[
\widetilde J_c(u;w)=g_c(u)+\frac{1}{2\mu_c}\|u-w\|^2
\]
est exactement la forme standard d'une enveloppe de Moreau, ce qui justifie la lecture proximale. Enfin, comparer les branches sur la même échelle évite qu'une rigidité quadratique plus forte domine artificiellement la sélection du contexte.
\end{proof}

\begin{definition}[Valeur réduite et enveloppe de Moreau]
Pour tout contexte \(c\) et toute entrée \(w\), on définit la valeur réduite
\[
\widetilde V_c(w):=\inf_{u\in\R^d}\widetilde J_c(u;w).
\]
On note également
\[
u_c^\star(w):=\argmin_{u\in\R^d}\widetilde J_c(u;w),
\qquad
e_{\mu_c g_c}(w):=\inf_{u\in\R^d}\left\{g_c(u)+\frac{1}{2\mu_c}\|u-w\|^2\right\}.
\]
\end{definition}

\begin{proposition}[Enveloppe de Moreau et gradient]
Pour tout contexte \(c\), la fonction \(g_c\) est propre, convexe et semi-continue inférieurement. On a
\[
u_c^\star(w)=\prox_{\mu_c g_c}(w),
\qquad
\widetilde V_c(w)=e_{\mu_c g_c}(w).
\]
De plus, \(e_{\mu_c g_c}\) est de classe \(C^1\) sur \(\R^d\) et
\[
\nabla e_{\mu_c g_c}(w)=\frac{1}{\mu_c}\bigl(w-\prox_{\mu_c g_c}(w)\bigr)
=\frac{1}{\mu_c}\bigl(w-u_c^\star(w)\bigr).
\]
Par conséquent,
\[
\nabla \widetilde V_c(w)=\frac{1}{\mu_c}\bigl(w-u_c^\star(w)\bigr).
\]
\end{proposition}

\begin{proof}
La convexité et la semi-continuité inférieure de \(g_c\) proviennent de celles de \(\dist(\cdot,\Mcal_c)^2\) et de \(\|\cdot\|^2\). Le résultat classique sur les enveloppes de Moreau d'une fonction propre, convexe et s.c.i. implique la différentiabilité \(C^1\) de \(e_{\mu_c g_c}\) et la formule du gradient. Comme \(\widetilde J_c(u;w)=g_c(u)+\frac{1}{2\mu_c}\|u-w\|^2\), les minimiseurs coïncident.
\end{proof}

\begin{remark}
Ces identités sont des résultats standards de l'analyse convexe ; elles ne sont invoquées que pour rendre le manuscrit autonome \cite{Moreau1965,Rockafellar1970,BauschkeCombettes2017,RockafellarWets1998}.
\end{remark}

% =================================================
\section{Routage exact, géométrie des cellules et régularité}
% =================================================

\begin{definition}[Fonctionnelle de routage]
Pour une entrée \(w\in\R^d\), on introduit
\[
\mathcal E_{\mathrm{route}}(w;u,c):=\widetilde J_c(u;w),
\qquad (u,c)\in\R^d\times\Ccal.
\]
\end{definition}

\begin{theorem}[Lower envelope routing theorem]
Pour tout \(w\in\R^d\),
\[
F(w):=\min_{c\in\Ccal}\widetilde V_c(w)
=
\inf_{(u,c)\in\R^d\times\Ccal}\mathcal E_{\mathrm{route}}(w;u,c).
\]
De plus, tout minimiseur global \((\hat u(w),\hat c(w))\) vérifie
\[
\hat u(w)=u_{\hat c(w)}^\star(w),
\qquad
\hat c(w)\in\argmin_{c\in\Ccal}\widetilde V_c(w).
\]
La fonction \(F\) est continue et localement lipschitzienne. Elle est directionnellement différentiable et, pour tout \(h\in\R^d\),
\[
F'(w;h)=\min_{c\in\Acal(w)}\langle \nabla \widetilde V_c(w),h\rangle.
\]
Enfin,
\[
\partial^C F(w)=\conv\{\nabla \widetilde V_c(w): c\in\Acal(w)\}.
\]
Sur toute cellule stricte \(\Omega_c^{\mathrm{str}}\), la fonction \(F\) est de classe \(C^1\) et
\[
\nabla F(w)=\nabla \widetilde V_c(w).
\]
\end{theorem}

\begin{proof}
Comme \(\Ccal\) est fini,
\[
\inf_{(u,c)\in\R^d\times\Ccal}\mathcal E_{\mathrm{route}}(w;u,c)
=
\min_{c\in\Ccal}\inf_{u\in\R^d}\widetilde J_c(u;w)
=
\min_{c\in\Ccal}\widetilde V_c(w).
\]
Le minimiseur en \(u\) est unique pour chaque contexte \(c\), d'où la première partie. Comme \(F\) est le minimum fini de fonctions \(C^1\), il est continu, localement lipschitzien et directionnellement différentiable. Les formules de dérivée directionnelle et de sous-différentiel de Clarke sont les règles standards du minimum fini de fonctions \(C^1\) localement lipschitziennes \cite{Clarke1983}. Sur une cellule stricte, une seule branche est active, donc \(F\) y coïncide avec une fonction \(C^1\) unique.
\end{proof}

\begin{proposition}[Correspondance active et version finie de Berge]
La correspondance active \(\Acal:\R^d\rightrightarrows\Ccal\) a des valeurs non vides et finies. Elle est semi-continue supérieurement. En particulier, sur toute cellule de dominance stricte, elle est localement constante. Ce point est la version finie du théorème du maximum de Berge \cite{Berge1963}.
\end{proposition}

\begin{proof}
Pour chaque \(c\),
\[
\{w\in\R^d:\ c\in\Acal(w)\}
=
\bigcap_{j\in\Ccal}\{w\in\R^d:\widetilde V_c(w)\le \widetilde V_j(w)\},
\]
intersection finie d'ensembles fermés, donc ensemble fermé. Comme \(\Ccal\) est fini, la correspondance est à valeurs finies, non vides et fermées ; elle est donc semi-continue supérieurement. Sur une cellule stricte, les inégalités sont strictes et stables par continuité, d'où la constance locale.
\end{proof}

\begin{definition}[Cellules de contexte]
On définit
\[
\Omega_c^{\mathrm{str}}:=\{w\in\R^d:\ \widetilde V_c(w)<\widetilde V_j(w)\ \forall j\neq c\},
\]
\[
\Omega_c:=\{w\in\R^d:\ c\in\argmin_{j\in\Ccal}\widetilde V_j(w)\}
=\bigcap_{j\neq c}\{w\in\R^d:\widetilde V_c(w)\le \widetilde V_j(w)\},
\]
et
\[
\Bcal:=\{w\in\R^d:\ \#\Acal(w)\ge 2\}.
\]
\end{definition}

\begin{proposition}[Partition de l'espace latent et frontières]
Pour chaque \(c\in\Ccal\), \(\Omega_c^{\mathrm{str}}\) est ouvert et \(\Omega_c\) est fermé. Les cellules strictes sont deux à deux disjointes, et
\[
\R^d=\left(\bigsqcup_{c\in\Ccal}\Omega_c^{\mathrm{str}}\right)\sqcup \Bcal.
\]
Pour chaque paire \(c\neq j\), en posant
\[
\Gamma_{c,j}(w):=\widetilde V_c(w)-\widetilde V_j(w),
\qquad
\Sigma_{c,j}:=\{w\in\R^d:\Gamma_{c,j}(w)=0\},
\]
on a que \(\Sigma_{c,j}\) est fermé. Si \(\nabla\Gamma_{c,j}(w_0)\neq 0\), alors \(\Sigma_{c,j}\) est localement une hypersurface \(C^1\) au voisinage de \(w_0\).
\end{proposition}

\begin{proof}
Chaque \(\widetilde V_c-\widetilde V_j\) est continue. Donc \(\Omega_c^{\mathrm{str}}\), défini par des inégalités strictes de fonctions continues, est ouvert, et \(\Omega_c\), défini par des inégalités larges, est fermé. Les cellules strictes sont disjointes par contradiction immédiate. La fermeture de \(\Sigma_{c,j}\) découle de la continuité de \(\Gamma_{c,j}\). Si \(\nabla\Gamma_{c,j}(w_0)\neq 0\), le théorème des fonctions implicites s'applique à \(\Gamma_{c,j}(w)=0\).
\end{proof}

\begin{corollary}[Stabilité sous marge stricte]
Supposons qu'en un point \(w\in\R^d\), le minimiseur de \(c\mapsto\widetilde V_c(w)\) soit unique, noté \(c^\star\), et qu'il existe une marge
\[
\delta(w):=\min_{c\neq c^\star}\bigl(\widetilde V_c(w)-\widetilde V_{c^\star}(w)\bigr)>0.
\]
Alors il existe un voisinage \(U\) de \(w\) tel que tout \(z\in U\) admet le même contexte routé \(c^\star\).
\end{corollary}

\begin{proof}
Par continuité finie des fonctions \(\widetilde V_c\), on choisit \(U\) tel que \(|\widetilde V_c(z)-\widetilde V_c(w)|<\delta(w)/3\) pour tous \(c\) et tout \(z\in U\). L'argument standard de comparaison des branches donne alors la stabilité locale du contexte optimal.
\end{proof}

% =================================================
\section{Réalisation algorithmique et profil de calcul}
% =================================================

Le routage exact peut être coûteux, car il exige la résolution complète de chaque problème contextuel. On introduit donc un certificat de marge et un schéma de pré-roulage court pour obtenir un routage calculable.

\begin{proposition}[Profil de calcul]
Si le coût d'une résolution locale du contexte \(c\) à précision \(\varepsilon\) est noté \(T_c(\varepsilon)\), alors un routage exact par comparaison exhaustive coûte
\[
C_{\mathrm{exact}}(\varepsilon)=\sum_{c\in\Ccal}T_c(\varepsilon).
\]
Un routage appris par gate a un coût d'inférence approximatif
\[
C_{\mathrm{gate}}\approx T_{\mathrm{gate}}+T_{c^\star},
\]
mais il exige une phase d'apprentissage distincte et n'offre pas, en lui-même, de certificat de marge sur le contexte choisi. Le routage variationnel exact est donc plus coûteux à l'inférence, mais il est certifiable au sens où une marge calculable suffit à valider la décision.
\end{proposition}

\begin{proof}
Le premier terme suit de la comparaison explicite de tous les contextes. Le second profil reflète le schéma classique gate-plus-expert : un gate est évalué puis un expert est exécuté. La différence conceptuelle est que le gate appris doit être entraîné et peut se tromper hors distribution, alors que le routage variationnel compare directement les valeurs locales et peut être validé par marge.
\end{proof}

\begin{definition}[Pré-roulage court]
Pour un contexte \(c\), on pose
\[
u_c^{(0)}(w)=w,
\qquad
u_c^{(n+1)}(w)=\Tcal_{c,\eta}(u_c^{(n)}(w);w),
\qquad n\ge 0.
\]
\end{definition}

\begin{definition}[Valeur de pré-roulage]
Pour un entier \(N\ge 1\), on définit
\[
\widetilde V_c^{(N)}(w):=\widetilde J_c(u_c^{(N)}(w);w).
\]
\end{definition}

\begin{assumption}[Hypothèse de contraction]
Pour un contexte \(c\) fixé, on suppose que l'opérateur
\[
\Tcal_{c,\eta}(z;w):=\argmin_{u\in\R^d}\left\{\widetilde J_c(u;w)+\frac{1}{2\eta}\|u-z\|^2\right\}
\]
est contractant sur \(\R^d\) avec facteur \(\rho_c\in(0,1)\), pour tout \(w\) dans le domaine considéré.
\end{assumption}

\begin{proposition}[Convergence géométrique et approximation de la valeur réduite]
Sous l'hypothèse de contraction, la suite \(u_c^{(N)}(w)\) converge vers \(u_c^\star(w)\) et
\[
\|u_c^{(N)}(w)-u_c^\star(w)\|\le \rho_c^N\|u_c^{(0)}(w)-u_c^\star(w)\|.
\]
De plus, si \(\widetilde J_c(\cdot;w)\) est Lipschitz sur un compact contenant l'orbite \(\{u_c^{(N)}(w)\}_{N\ge 0}\cup\{u_c^\star(w)\}\), alors il existe \(L_{c,w}<\infty\) tel que
\[
\bigl|\widetilde V_c^{(N)}(w)-\widetilde V_c(w)\bigr|\le L_{c,w}\rho_c^N\|u_c^{(0)}(w)-u_c^\star(w)\|.
\]
\end{proposition}

\begin{proof}
La convergence géométrique des itérés est la conclusion standard du théorème du point fixe de Banach pour une contraction de facteur \(\rho_c\). Comme \(\widetilde J_c\) est continue et \(C^1\), elle est Lipschitz sur tout compact. En évaluant la différence des valeurs sur le compact contenant l'orbite et le point fixe, on obtient l'inégalité annoncée.
\end{proof}

\begin{theorem}[Stabilisation du routage approché]
Supposons que le minimiseur de \(c\mapsto\widetilde V_c(w)\) soit unique, noté \(c^\star(w)\). Alors il existe \(N_0\ge 1\) tel que, pour tout \(N\ge N_0\),
\[
\hat c_N(w)=c^\star(w),
\qquad
\hat c_N(w)\in\argmin_{c\in\Ccal}\widetilde V_c^{(N)}(w).
\]
\end{theorem}

\begin{proof}
Comme \(\Ccal\) est fini, \(\argmin_{c\in\Ccal}\widetilde V_c^{(N)}(w)\) est non vide pour tout \(N\). En outre, pour chaque \(c\), \(\widetilde V_c^{(N)}(w)\to \widetilde V_c(w)\). La convergence est donc uniforme sur \(\Ccal\). Posons
\[
\delta:=\min_{c\neq c^\star(w)}\bigl(\widetilde V_c(w)-\widetilde V_{c^\star(w)}(w)\bigr)>0.
\]
Il existe alors \(N_0\ge 1\) tel que, pour tout \(N\ge N_0\) et tout \(c\in\Ccal\),
\[
\bigl|\widetilde V_c^{(N)}(w)-\widetilde V_c(w)\bigr|<\frac{\delta}{3}.
\]
Les mêmes comparaisons que dans la stabilité sous marge stricte donnent alors \(\hat c_N(w)=c^\star(w)\) pour tout \(N\ge N_0\).
\end{proof}

\begin{corollary}[Routage approché par certificat court]
On peut définir un routage approché par
\[
\hat c_N(w)\in\argmin_{c\in\Ccal}\widetilde V_c^{(N)}(w).
\]
Lorsque \(N\) est suffisamment grand, \(\hat c_N(w)\) coïncide avec le routage exact sous l'hypothèse d'unicité du minimiseur.
\end{corollary}

\begin{quote}
\small
\textbf{Algorithme de routage calculable.}

Entrée : \(w\), ensemble fini \(\Ccal\), tolérance \(\varepsilon\).

1. Pour chaque contexte \(c\in\Ccal\), résoudre le problème local et calculer une approximation \(\widetilde V_c^{(N)}(w)\).

2. Choisir \(c^\star\in\argmin_{c\in\Ccal}\widetilde V_c^{(N)}(w)\).

3. Si une marge \(\delta>0\) est certifiée, arrêter dès que
\[
\bigl|\widetilde V_c^{(N)}(w)-\widetilde V_{c^\star}^{(N)}(w)\bigr|<\delta/3
\qquad\text{pour tout }c\neq c^\star.
\]

Sortie : contexte routé \(c^\star\) et minimiseur local \(u_{c^\star}^\star(w)\).
\end{quote}

\begin{remark}
La procédure précédente n'est pas une heuristique cachée : elle n'utilise que des valeurs locales, une marge calculable et un critère d'arrêt explicite.
\end{remark}

% =================================================
\section{Routage souple comme relaxation entropique}
% =================================================

Le routage dur est discret. Pour obtenir une relaxation convexe, on passe au simplexe des contextes.

\begin{definition}[Simplexe des contextes]
On note
\[
\Delta_K:=\left\{p\in\R^K:\ p_c\ge 0,\ \sum_{c=1}^K p_c=1\right\},
\qquad
\Delta_K^\circ:=\{p\in\Delta_K:\ p_c>0\ \text{pour tout }c\}.
\]
\end{definition}

\begin{definition}[Fonctionnelle de routage souple]
Pour \(\beta>0\), on définit
\[
\Rcal_\beta(w,p)
=
\sum_{c=1}^K p_c\,\widetilde V_c(w)
+\frac{1}{\beta}\sum_{c=1}^K p_c\log p_c,
\qquad p\in\Delta_K,
\]
avec la convention \(0\log 0=0\).
\end{definition}

\begin{proposition}[Existence et unicité du routage souple]
Pour tout \(w\), la fonctionnelle \(p\mapsto \Rcal_\beta(w,p)\) admet un unique minimiseur \(p^\star_\beta(w)\) dans \(\Delta_K\).
\end{proposition}

\begin{proof}
La fonction \(p\mapsto \sum_c p_c\log p_c\) est strictement convexe sur \(\Delta_K\) et \(p\mapsto \sum_c p_c\,\widetilde V_c(w)\) est affine. La somme est donc strictement convexe sur le compact \(\Delta_K\), d'où l'existence et l'unicité du minimiseur.

Montrons maintenant que le minimiseur appartient à l'intérieur \(\Delta_K^\circ\). Supposons au contraire que \(p^\star_i=0\) pour un certain indice \(i\). Comme \(\sum_c p^\star_c=1\), il existe un indice \(j\) tel que \(p^\star_j>0\). Pour \(t\in(0,p^\star_j)\), posons \(p(t)=p^\star+t(e_i-e_j)\in\Delta_K\). On obtient une fonction unidimensionnelle dont la dérivée à droite en \(0\) est \(-\infty\), car \(\log t\to-\infty\) lorsque \(t\downarrow 0\). Cela contredit la minimalité de \(p^\star\). Ainsi, tout minimiseur appartient à \(\Delta_K^\circ\).
\end{proof}

\begin{theorem}[Forme explicite du routage souple]
Le minimiseur unique de \(\Rcal_\beta(w,\cdot)\) est donné par
\[
p_c^\star(w)=\frac{\exp\!\bigl(-\beta\,\widetilde V_c(w)\bigr)}
{\sum_{j=1}^K \exp\!\bigl(-\beta\,\widetilde V_j(w)\bigr)}.
\]
\end{theorem}

\begin{proof}
Comme le minimiseur appartient à \(\Delta_K^\circ\), on peut appliquer la méthode des multiplicateurs de Lagrange au problème
\[
\min_{p\in\Delta_K^\circ}\ \sum_{c=1}^K p_c\,\widetilde V_c(w)+\frac{1}{\beta}\sum_{c=1}^K p_c\log p_c.
\]
On obtient, pour chaque \(c\),
\[
\widetilde V_c(w)+\frac{1}{\beta}(1+\log p_c)+\lambda=0,
\]
donc
\[
p_c=\exp(-1-\beta\lambda)\,\exp\!\bigl(-\beta\widetilde V_c(w)\bigr).
\]
La contrainte \(\sum_c p_c=1\) donne la formule annoncée. Le résultat est la forme usuelle du softmin entropique \cite{BauschkeCombettes2017,RockafellarWets1998}.
\end{proof}

\begin{corollary}[Valeur du problème de routage et limite dure]
On a l'identité
\[
\inf_{p\in\Delta_K}\Rcal_\beta(w,p)
=
-\frac{1}{\beta}\log\!\left(\sum_{c=1}^K \exp\!\bigl(-\beta\,\widetilde V_c(w)\bigr)\right).
\]
De plus,
\[
F(w)\le -\frac{1}{\beta}\log\!\left(\sum_{c=1}^K e^{-\beta \widetilde V_c(w)}\right)\le F(w)+\frac{\log K}{\beta}.
\]
Enfin, si le minimiseur de \(c\mapsto\widetilde V_c(w)\) est unique, noté \(c^\star(w)\), alors
\[
p^\star_\beta(w)\xrightarrow[\beta\to+\infty]{}e_{c^\star(w)}.
\]
\end{corollary}

\begin{proof}
C'est la valeur classique de l'entropie log-sum-exp obtenue en injectant la solution précédente dans \(\Rcal_\beta\). La borne sandwich est la propriété standard du \(\log\)-sum-exp. Pour la limite dure, on compare chaque poids non optimal au poids optimal ; les différences strictement positives dans les valeurs réduites sont écrasées exponentiellement lorsque \(\beta\to+\infty\).
\end{proof}

\begin{remark}
Le routage souple est une relaxation entropique du problème discret. Il n'abolit pas le routage exact ; il le lisse de manière stable.
\end{remark}

% =================================================
\section{Pont avec PCH-ESM : décomposition spectrale et choix du bloc}
% =================================================

Dans PCH-ESM, une décomposition orthogonale et déterministe fournit une famille finie de sous-espaces spécialisés. Ce sont des contextes naturels pour le routage variationnel.

\begin{assumption}[Décomposition spectrale]
On suppose qu'un latent global \(z\in\Hh\) admet une décomposition orthogonale
\[
\Hh=\mathfrak S\oplus\mathfrak S^\perp,
\]
où \(\mathfrak S\) est la somme directe finie des sous-espaces spécialisés \(\Hh_1,\dots,\Hh_n\). On note \(\Pi_{\mathfrak S}\) le projecteur orthogonal sur \(\mathfrak S\).
\end{assumption}

\begin{definition}[Contextes PCH-ESM]
On note \(\Ccal_{\mathrm{ESM}}=\{1,\dots,n\}\) l'ensemble des blocs spécialisés induits par la décomposition spectrale. À chaque bloc \(i\) on associe une valeur réduite \(\widetilde V_i^{\mathrm{ESM}}(w)\) calculée par le problème local du bloc correspondant.
\end{definition}

\begin{proposition}[Instantiation du routage sur PCH-ESM]
Si les contextes admissibles sont pris parmi les blocs spécialisés de PCH-ESM, alors le bloc routé est
\[
i^\star(w)\in\argmin_{i\in\Ccal_{\mathrm{ESM}}}\widetilde V_i^{\mathrm{ESM}}(w).
\]
Autrement dit, la décomposition spectrale fournit une famille de contextes admissibles, et le routage variationnel sélectionne le bloc de plus faible valeur réduite.
\end{proposition}

\begin{proof}
C'est l'application directe du théorème de routage exact à la famille finie des blocs spécialisés.
\end{proof}

\begin{proposition}[Amélioration par spécialisation orthogonale]
Supposons que la cible vraie \(\theta^\star\) appartienne à \(\mathfrak S\), et qu'une estimation dense \(\hat\theta_{\mathrm{dense}}\in\Hh\) soit projetée par \(\hat\theta_{\mathrm{spec}}:=\Pi_{\mathfrak S}\hat\theta_{\mathrm{dense}}\). Alors
\[
\|\hat\theta_{\mathrm{dense}}-\theta^\star\|^2
=
\|\hat\theta_{\mathrm{spec}}-\theta^\star\|^2
+
\|\Pi_{\mathfrak S^\perp}\hat\theta_{\mathrm{dense}}\|^2.
\]
En particulier, si \(\Pi_{\mathfrak S^\perp}\hat\theta_{\mathrm{dense}}\neq 0\), la spécialisation réduit strictement l'erreur quadratique.
\end{proposition}

\begin{proof}
Comme \(\theta^\star\in\mathfrak S\),
\[
\hat\theta_{\mathrm{dense}}-\theta^\star
=
\bigl(\Pi_{\mathfrak S}\hat\theta_{\mathrm{dense}}-\theta^\star\bigr)+\Pi_{\mathfrak S^\perp}\hat\theta_{\mathrm{dense}}.
\]
Les deux termes sont orthogonaux. Le résultat suit par Pythagore.
\end{proof}

\begin{figure}[ht]
\centering
\begin{tikzpicture}[node distance=1.6cm, >=Latex, every node/.style={align=center}]
\node (pch) [draw, rounded corners, minimum width=3.2cm, minimum height=0.95cm] {PCH\\valeurs locales};
\node (esm) [draw, rounded corners, right=1.8cm of pch, minimum width=3.5cm, minimum height=0.95cm] {PCH-ESM\\blocs spécialisés};
\node (vals) [draw, rounded corners, right=1.8cm of esm, minimum width=3.4cm, minimum height=0.95cm] {Valeurs réduites\\par bloc};
\node (sel) [draw, rounded corners, below=1.2cm of vals, minimum width=3.4cm, minimum height=0.95cm] {Sélection\\\(\argmin\)};
\node (out) [draw, rounded corners, left=1.8cm of sel, minimum width=3.4cm, minimum height=0.95cm] {Contexte choisi\\ou bloc expert};

\draw[->] (pch) -- (esm);
\draw[->] (esm) -- (vals);
\draw[->] (vals) -- (sel);
\draw[->] (sel) -- (out);
\end{tikzpicture}
\caption{Chaîne conceptuelle PCH \(\rightarrow\) PCH-ESM \(\rightarrow\) routage.}
\end{figure}

\begin{remark}
Cette compatibilité clarifie le rôle de PCH-ESM : il ne remplace pas le routage, il fournit une famille de branches spécialisées sur lesquelles le routage variationnel opère.
\end{remark}

% =================================================
\section{Théorème structurel principal}
% =================================================

\begin{theorem}[Structure variationnelle globale et instantiation PCH-ESM]
Sous les hypothèses du document :
\begin{enumerate}[leftmargin=2em]
\item pour chaque contexte \(c\), le problème local admet un unique minimiseur \(u_c^\star(w)=\prox_{\mu_c g_c}(w)\) et \(\widetilde V_c(w)=e_{\mu_c g_c}(w)\) ;
\item la fonction globale \(F(w)=\min_{c\in\Ccal}\widetilde V_c(w)\) est continue, localement lipschitzienne, directionnellement différentiable, et de classe \(C^1\) sur toute cellule stricte ;
\item l'espace latent se décompose en cellules strictes \(\Omega_c^{\mathrm{str}}\), en cellules faibles \(\Omega_c\), et en une strate de bord \(\Bcal\) ;
\item la correspondance active est semi-continue supérieurement et la régularité de Clarke est donnée par l'enveloppe convexe des gradients actifs ;
\item le routage souple est la régularisation entropique du critère des valeurs réduites, avec minimiseur de Gibbs ;
\item si les contextes admissibles proviennent des blocs spécialisés de PCH-ESM, alors le bloc routé est exactement le bloc minimisant la valeur réduite correspondante ;
\item la spécialisation orthogonale améliore strictement l'erreur quadratique dès lors qu'une composante de l'estimation dense vit hors du sous-espace spécialisé.
\end{enumerate}
\end{theorem}

\begin{proof}
Les points \(1\) à \(5\) sont établis dans les sections précédentes. Le point \(6\) est l'instantiation du principe de routage à la famille de blocs spécialisés fournie par PCH-ESM. Le point \(7\) résulte de la décomposition orthogonale et de l'identité de Pythagore.
\end{proof}

% =================================================
\section{Exemple illustratif en dimension un}
% =================================================

\begin{example}[Deux contextes symétriques]
On prend \(d=1\), \(\Mcal_1=\{-a\}\), \(\Mcal_2=\{a\}\) avec \(a>0\), et on suppose les mêmes paramètres pour les deux contextes :
\[
\lambda_{\mathrm{ctx},1}=\lambda_{\mathrm{ctx},2}=\lambda>0,
\qquad
\lambda_{\mathrm{ridge},1}=\lambda_{\mathrm{ridge},2}=\rho\ge 0.
\]
Alors \(\mu_1=\mu_2=1+2\rho\), et une complétion du carré donne
\[
u_1^\star(w)=\frac{w-2\lambda a}{1+2\lambda+2\rho},
\qquad
u_2^\star(w)=\frac{w+2\lambda a}{1+2\lambda+2\rho}.
\]
De plus,
\[
\widetilde V_1(w)-\widetilde V_2(w)
=
\frac{4a\lambda}{(1+2\rho)(1+2\lambda+2\rho)}\,w.
\]
Par conséquent,
\[
\Omega_1^{\mathrm{str}}=(-\infty,0),
\qquad
\Omega_2^{\mathrm{str}}=(0,\infty),
\qquad
\Bcal=\{0\}.
\]
Le routage exact se lit donc immédiatement : le contexte \(1\) est choisi pour \(w<0\), le contexte \(2\) pour \(w>0\), et la frontière de décision est ici réduite au point \(0\).
\end{example}

\begin{figure}[ht]
\centering
\begin{tikzpicture}[scale=1]
\draw[->] (-3.2,0) -- (3.2,0) node[right] {$x$};
\draw[->] (0,-0.2) -- (0,2.2) node[above] {$y$};
\draw[thick] (-2.8,0.3) .. controls (-1.8,1.6) and (-0.9,1.9) .. (0,2.05);
\draw[thick] (-2.8,1.9) .. controls (-1.7,1.3) and (-0.8,0.8) .. (0,0.45);
\draw[thick] (0,0.45) .. controls (0.8,0.8) and (1.7,1.3) .. (2.8,1.9);
\draw[dashed] (0,0) -- (0,2.1);
\node at (-1.85,1.0) {$V_1$};
\node at (1.85,1.0) {$V_2$};
\node at (0.55,0.18) {$\Bcal$};
\end{tikzpicture}
\caption{Lecture schématique de deux branches et de leur frontière.}
\end{figure}

\begin{remark}
Cet exemple montre que les cellules ne sont pas un objet abstrait : elles décrivent une partition géométrique explicite de l'espace latent.
\end{remark}

% =================================================
\section{Relation à la littérature et positionnement}
% =================================================

Le cadre ici présenté se place à l'interface de l'analyse convexe, de la théorie des fonctions valeur, de la programmation disjonctive et des modèles de type \emph{mixture of experts}. Les références classiques pertinentes sont celles de Moreau, Rockafellar, Berge, Clarke, ainsi que Bauschke--Combettes et Rockafellar--Wets \cite{Moreau1965,Rockafellar1970,Berge1963,Clarke1983,BauschkeCombettes2017,RockafellarWets1998}.

La nouveauté du manuscrit n'est pas d'inventer de nouveaux outils de convexité ; elle est de faire du routage un problème variationnel à part entière. Les méthodes classiques de routage utilisent un score appris, un classifieur ou un gate séparé. Ici, le contexte est sélectionné comme minimiseur d'une famille de valeurs réduites issues directement de problèmes locaux. Le routage n'est donc plus un module externe : il devient une conséquence de la structure énergétique.

En termes pratiques, le routage appris est souvent moins coûteux à l'inférence une fois le gate entraîné, mais il introduit une phase d'apprentissage séparée, une dépendance aux données d'entraînement du gate, et une opacité de décision. Le routage variationnel est plus coûteux si l'on résout exhaustivement toutes les branches, mais il est calculable, il s'inscrit directement dans PCH, et il peut être accéléré par un pré-roulage court avec marge. Ce compromis est la formulation scientifique du problème.

Autrement dit, la quantité fondamentale n'est pas un score appris séparément, mais l'objet
\[
\min_{c\in\Ccal} e_{\mu_c g_c}(w).
\]
C'est cette enveloppe inférieure qui porte l'information de décision, et non une sortie de classifieur externe.

\begin{remark}
Cette lecture permet de séparer proprement ce qui relève d'un rappel classique et ce qui relève de la contribution propre du manuscrit : la sélection de contexte par minimisation des valeurs réduites et l'interprétation géométrique qui en découle.
\end{remark}

% =================================================
\section{Conclusion}
% =================================================

Le noyau PCH n'est pas le problème à résoudre ; il est déjà la brique mathématique centrale. Le routage est la continuité naturelle de PCH : le contexte sélectionné est celui qui minimise la valeur locale réduite.

La formulation canonique est
\[
(\hat u(w),\hat c(w))\in\argmin_{(u,c)\in\R^d\times\Ccal}\mathcal E_{\mathrm{route}}(w;u,c),
\]
équivalente à
\[
\hat c(w)\in\argmin_{c\in\Ccal}\widetilde V_c(w),
\qquad
\widetilde V_c(w)=e_{\mu_c g_c}(w).
\]
La version souple est donnée par
\[
p^\star_\beta(w)=\argmin_{p\in\Delta_K}\Rcal_\beta(w,p).
\]

Cette formulation conserve strictement le noyau PCH, définit explicitement le critère de sélection du contexte, donne une écriture exacte du routage comme problème sur \(\R^d\times\Ccal\), admet une relaxation convexe sur le simplexe, relie le pré-roulage à la convergence vers la valeur réduite, et fournit une interprétation variationnelle propre du choix de contexte. Le pont avec PCH-ESM est également explicite : la décomposition spectrale fournit une famille finie de blocs spécialisés, et le même principe de minimisation sélectionne le bloc de plus faible valeur réduite.

\begin{thebibliography}{99}

\bibitem{Moreau1965}
J.-J. Moreau,
\newblock Proximité et dualité dans un espace hilbertien.
\newblock \emph{Bull. Soc. Math. France}, 93:273--299, 1965.

\bibitem{Rockafellar1970}
R. T. Rockafellar,
\newblock \emph{Convex Analysis}.
\newblock Princeton University Press, 1970.

\bibitem{Berge1963}
C. Berge,
\newblock \emph{Topological Spaces: Including a Treatment of Multi-Valued Functions, Vector Spaces and Convexity}.
\newblock Oliver \& Boyd, 1963.

\bibitem{Clarke1983}
F. H. Clarke,
\newblock \emph{Optimization and Nonsmooth Analysis}.
\newblock Wiley, 1983.

\bibitem{BauschkeCombettes2017}
H. H. Bauschke and P. L. Combettes,
\newblock \emph{Convex Analysis and Monotone Operator Theory in Hilbert Spaces}.
\newblock 2nd ed., Springer, 2017.

\bibitem{RockafellarWets1998}
R. T. Rockafellar and R. J.-B. Wets,
\newblock \emph{Variational Analysis}.
\newblock Springer, 1998.

\end{thebibliography}

\end{document}
